\section*{Bab \ref{ch:b2:blood-pressure} --- Pengaturan Tekanan Darah dan Olahraga}

\begin{solution}{exo:b2:blood-pressure:1}
$\bar P - P_{\text{vena}} = Q\,R = f\,V_s\,R$. Laju denyut \omterm{def:b2:heart:anatomy}{jantungnya}
(\omterm{prop:b2:heart:conduction}{nodus sinoatriumnya}, di bawah \omterm{prop:b2:heart:control}{saraf vagus} dan simpatisnya); \omterm{prop:b2:heart:cycle}{volume
sekuncupnya} (\omterm{def:b2:heart:anatomy}{biliknya}: pengisiannya, yang ditetapkan \omterm{def:b2:blood-circulation:vessels}{venanya}, dan daya
kerutnya, yang ditetapkan saraf simpatisnya); hambatannya (otot polos
\omterm{def:b2:blood-circulation:vessels}{arteriolnya}, di bawah saraf simpatis, hormon, dan metabolit
setempatnya).
\end{solution}

\begin{solution}{exo:b2:blood-pressure:2}
Pengindranya: reseptor regangan di dalam sinus karotis dan lengkung
aortanya. Pusatnya: pusat kardiovaskular medulanya. Efektornya: \omterm{prop:b2:heart:control}{saraf
vagus} (laju) dan saraf simpatis (laju, daya kerut, tonus \omterm{def:b2:blood-circulation:vessels}{arteriol} dan
\omterm{def:b2:blood-circulation:vessels}{vena}). Umpan balik negatif, dengan tundaan kira-kira sedetik. Tanpa
pengindranya tekanan rerata sepanjang sehari tetap normal tetapi
tekanan dari menit ke menit mengayun liar mengikuti tiap sikap tubuh
dan perasaan.
\end{solution}

\begin{solution}{exo:b2:blood-pressure:3}
Ginjalnya (sel \omterm{def:b2:blood-circulation:vessels}{arteriolnya}) melepaskan renin ketika pengalirannya atau
antaran natriumnya turun atau saraf simpatisnya memicu; renin memotong
angiotensinogen (hati) menjadi angiotensin I; enzim parunya mengubahnya
menjadi angiotensin II, yang menyempitkan \omterm{def:b2:blood-circulation:vessels}{arteriolnya} dan merangsang
korteks adrenalnya untuk menyekresikan aldosteron; aldosteron membuat
ginjalnya menahan natrium dan air, sehingga menaikkan volume \omterm{def:b2:blood-circulation:blood}{darahnya}.
\end{solution}

\begin{solution}{exo:b2:blood-pressure:4}
Pelebaran metabolis \omterm{def:b2:blood-circulation:vessels}{arteriol} ototnya sendiri oleh metabolitnya; \omterm{thm:b2:heart:starling}{curah
jantung} lima kali lipat yang digerakkan saraf simpatis dan pompa
ototnya; penyempitan simpatis pada \omterm{def:b2:blood-circulation:vessels}{arteriol} usus, ginjal, dan otot yang
beristirahat, yang menyerahkan alirannya.
\end{solution}

\begin{solution}{exo:b2:blood-pressure:5}
$83\times 1.2 = \qty{100}{mmHg}$; $333\times 0.35 = \qty{117}{mmHg}$.
Hambatannya turun sebesar $1.2/0.35 = 3.4$: jari-jarinya naik sebesar
$3.4^{1/4} =
1.36$.
\end{solution}

\begin{solution}{exo:b2:blood-pressure:6}
$30/(1 + G)$: 10, 6, dan \qty{3.3}{mmHg}. Dengan separuh batinya
penurunan sisanya melipat dua: ia pusing ketika berdiri, melihat
kelabu, dan mungkin pingsan.
\end{solution}

\begin{solution}{exo:b2:blood-pressure:7}
$\rho g h$: $1060\times 9.81\times 0.4 = \qty{4160}{Pa} =
\qty{31}{mmHg}$; $1060\times 9.81\times 1.2 = \qty{12500}{Pa} =
\qty{94}{mmHg}$. Otaknya $95 - 31 = \qty{64}{mmHg}$; pergelangan
kakinya $95 + 94 =
\qty{189}{mmHg}$. Kepala seekor jerapah pada
\qty{2.5}{m} berongkos \qty{195}{mmHg}: ia memerlukan tekanan rerata
sekitar \qty{250}{mmHg} di \omterm{def:b2:heart:anatomy}{jantungnya}, serta dinding pembuluh yang
tebal dan kulit yang ketat pada tungkainya.
\end{solution}

\begin{solution}{exo:b2:blood-pressure:8}
$Q = 280/(190 - 130) = \qty{4.7}{L/min}$. Atletnya: $5000/(200 - 30) =
\qty{29}{L/min}$; \omterm{prop:b2:heart:cycle}{volume sekuncupnya} $\num{29400}/190 = \qty{155}{mL}$.
\end{solution}

\begin{solution}{exo:b2:blood-pressure:9}
$35/5 = \qty{7}{mmHg}$ turun: kira-kira \qty{88}{mmHg}. Memulihkan
volumenya: penyerapan kembali cairan antarsel ke dalam \omterm{def:b2:blood-circulation:vessels}{kapilernya}
(beberapa menit sampai sejam); air seni yang berkurang di bawah \omterm{prop:b2:blood-pressure:raas}{hormon
antidiuretik} dan aldosteron (berjam-jam sampai berhari-hari); rasa haus
dan minum (berjam-jam); \omterm{def:b2:blood-circulation:blood}{sel darah merah} dari sumsumnya
(berminggu-minggu). Kulitnya pucat dan dingin karena refleksnya telah
menutup \omterm{def:b2:blood-circulation:vessels}{arteriolnya} demi menyimpan tekanannya bagi otak dan \omterm{def:b2:heart:anatomy}{jantungnya}.
\end{solution}

\begin{solution}{exo:b2:blood-pressure:10}
Otot pada $R = 0.5$: $95/0.5 = \qty{190}{mL/s}$, \qty{11.4}{L/min};
bersama \qty{4}{L/min} organ lainnya \omterm{def:b2:heart:anatomy}{jantungnya} akan memerlukan
\qty{15.4}{L/min}. Dengan $Q$ dipancangkan pada \qty{5}{L/min},
hambatan totalnya turun dari 1.14 menjadi $1/(1/1.14 - 1/5 + 1/0.5) =
\qty{0.37}{mmHg.s/mL}$ dan tekanannya menjadi \qty{31}{mmHg} ---
otaknya akan pingsan. Tubuhnya menaikkan keluarannya menuju keadaan
yang pertama lalu menyempitkan petak lainnya, sehingga tekanannya
bertahan dan ototnya mendapatkan alirannya.
\end{solution}

\begin{solution}{exo:b2:blood-pressure:11}
(a) Ginjal yang kurang teraliri itu membaca tekanan yang rendah,
melepaskan renin, lalu angiotensin dan aldosteron menaikkan hambatan
dan volumenya. (b) \omterm{def:b2:blood-pressure:baroreflex}{Barorefleksnya} menyangga perubahan tetapi menyetel
ulang dirinya kepada aras mana pun yang bertahan, dan tidak dapat
mengubah penahanan volume oleh ginjalnya. (c) Sumber reninnya, dan
organ yang menuntut tekanan yang lebih tinggi, telah lenyap. (d)
Menghalangi pengubahan menjadi angiotensin II: hambatan yang lebih
rendah, aldosteron yang lebih sedikit, tekanan yang lebih rendah.
\end{solution}

\begin{solution}{exo:b2:blood-pressure:12}
\omterm{def:b2:blood-pressure:baroreflex}{Barorefleksnya} bekerja dalam sedetik dengan bati kira-kira empat dan
titik acuan yang menghanyut menuju tekanan yang berlaku: ia meredam
lonjakan dan tidak dapat menahan sebuah aras. Ginjalnya bekerja selama
berjam-jam dan berhari-hari, dengan bati yang pada dasarnya tak
terhingga --- ia terus membuang atau menahan sampai tekanan yang
menjadi setelannya tercapai --- dan titik acuannya dipancangkan oleh
faalnya sendiri: ia menahan arasnya. Karena itu kenaikan yang menahun
berarti titik acuan ginjalnya telah berpindah, dan hanya obat yang
bekerja pada gelung ginjalnya atau pada hambatan yang ia tuntut yang
menurunkannya.
\end{solution}

\begin{solution}{pb:b2:blood-pressure:1}
\textbf{1.} $1/R = 1/7.6 + 1/22.8 + 1/4.1 + 1/5.2 + 1/5.7 + 1/19 +
1/28.5 = 0.875$: $R = \qty{1.14}{mmHg.s/mL}$; $\bar P/Q = 95/83.3 =
1.14$.
\textbf{2.} $\bar P/R_i$: otak 0.75, \omterm{def:b2:heart:anatomy}{jantung} 0.25, usus dan hati
1.39, ginjal 1.10, otot 1.0, kulit 0.30, lainnya \qty{0.20}{L/min};
bagiannya 15, 5, 28, 22, 20, 6, dan \qty{4}{\%}.
\textbf{3.} $5000/70 = \qty{71}{mL}$; $5\times 50 = \qty{250}{mL}$
oksigen tiap menit.
\textbf{4.} Otaknya $0.75/1.4 = \qty{0.54}{L/min/kg}$; ginjalnya $1.1/0.3
= \qty{3.7}{L/min/kg}$, yaitu tujuh kali otaknya dan lima puluh kali
rerata tubuhnya: ginjalnya dialiri untuk menapis, bukan untuk diberi makan.
\textbf{5.} Dengan satu tekanan yang dibagi seluruh organnya,
masing-masing menarik aliran yang ditetapkan hambatannya sendiri;
mengatur alirannya secara memusat akan membuat kelaparan organ mana pun
yang permintaannya berubah. Tekanan adalah mata uang bersamanya.
\textbf{6.} \omterm{def:b2:heart:anatomy}{Jantungnya}: $250\times 0.2\times 0.7 = \qty{35}{mL/min}$;
ototnya: $1000\times 0.2\times 0.25 = \qty{50}{mL/min}$. \omterm{def:b2:heart:anatomy}{Jantungnya}
sudah menyaripatikan sebagian besar dari yang lewat, sehingga ia hanya
dapat memperoleh oksigen lewat aliran yang lebih banyak, dan itulah
sebabnya \omterm{def:b2:blood-circulation:vessels}{arteriolnya} melebar begitu ia bekerja lebih keras.
\textbf{7.} $Q$ turun \qty{30}{\%} pada $R$ yang tetap: $\bar P$ turun
\qty{30}{\%}, \qty{28.5}{mmHg}, menjadi \qty{66.5}{mmHg}.
\textbf{8.} $28.5/5 = \qty{5.7}{mmHg}$: kira-kira \qty{89}{mmHg}.
\textbf{9.} $Q = 85\times 0.7\times 71 = \qty{4.2}{L/min} =
\qty{70}{mL/s}$; $R = 89.3/70 = \qty{1.27}{mmHg.s/mL}$, naik
\qty{11}{\%}.
\textbf{10.} $1060\times 9.81\times 0.4 = \qty{4160}{Pa} =
\qty{31}{mmHg}$; otaknya $66.5 - 31 = \qty{35}{mmHg}$ sebelum
refleksnya, $89 - 31 = \qty{58}{mmHg}$ sesudahnya.
\textbf{11.} $28.5/2 = \qty{14}{mmHg}$: \qty{81}{mmHg} di \omterm{def:b2:heart:anatomy}{jantungnya},
\qty{50}{mmHg} di otaknya --- pusing, penglihatan kelabu, hampir
pingsan tiap kali bangkit dari sebuah kursi.
\textbf{12.} Berbaring datar menghapus rugi hidrostatik
\qty{31}{mmHg} dan mengalirkan \omterm{def:b2:blood-circulation:blood}{darah} yang menggenang kembali ke
\omterm{def:b2:heart:anatomy}{jantungnya}; \omterm{prop:b2:heart:cycle}{volume sekuncup} dan tekanan otaknya pulih dalam beberapa
denyut.
\textbf{13.} $45/5 = \qty{9}{mmHg}$: \qty{86}{mmHg}.
\textbf{14.} $1/R = 0.875 - 1/19 - 1/4.1 + 1/38 + 1/8.2 = 0.727$: $R =
\qty{1.38}{mmHg.s/mL}$. Otaknya $86/7.6 = \qty{11.3}{mL/s} =
\qty{0.68}{L/min}$; kulitnya $86/38 = \qty{2.3}{mL/s} = \qty{0.14}{L/min}$;
ususnya $86/8.2 = \qty{10.5}{mL/s} = \qty{0.63}{L/min}$.
\textbf{15.} Dengan \omterm{def:b2:blood-circulation:blood}{darah} yang lebih sedikit dan \omterm{def:b2:blood-circulation:vessels}{arteriol} yang lebih
kencang, tekanan \omterm{def:b2:blood-circulation:vessels}{kapilernya} turun di bawah \omterm{thm:b2:blood-circulation:oncotic}{tekanan onkotiknya} sepanjang
seluruh \omterm{def:b2:blood-circulation:vessels}{kapilernya}, sehingga cairannya diserap kembali alih-alih
disaring; \omterm{def:b2:blood-circulation:blood}{plasmanya} mengencer dan \omterm{def:b2:blood-circulation:blood}{hematokritnya} turun.
\textbf{16.} Aldosteron, lewat renin dan angiotensin, yang dipicu oleh
pengaliran ginjalnya yang rendah dan dorongan simpatisnya; \omterm{prop:b2:blood-pressure:raas}{hormon
antidiuretik}, yang dipicu oleh penurunan volumenya dan kenaikan
kepekatan \omterm{def:b2:blood-circulation:blood}{plasmanya}.
\textbf{17.} $5\times 10^{12}$ sel pada $2.5\times 10^{6}$ tiap detik:
$2\times 10^{6}$ s, kira-kira tiga pekan.
\textbf{18.} Di bawah kira-kira \qty{60}{mmHg} kurva refleksnya datar:
pembuluhnya sudah sekencang mungkin dan \omterm{def:b2:heart:anatomy}{jantungnya} secepat mungkin,
batinya nol, dan tiap kehilangan selanjutnya menurunkan tekanannya
sebesar jumlah penuhnya; jaringan yang kurang teraliri lalu melepaskan
asam dan melebar, dan tekanannya runtuh.
\textbf{19.} Ototnya $110/0.33 = \qty{333}{mL/s} = \qty{20}{L/min}$;
kulitnya $110/5 = \qty{22}{mL/s} = \qty{1.3}{L/min}$.
\textbf{20.} Hambatan otaknya $8.8$ (aliran 0.75 pada 110), \omterm{def:b2:heart:anatomy}{jantungnya} $6.6$
(aliran 1.0): $1/R = 1/8.8 + 1/6.6 + 1/8.2 + 1/10.4 + 1/0.33 + 1/5 +
1/28.5 = 3.75$: $R = \qty{0.27}{mmHg.s/mL}$; $Q = 110/0.27 =
\qty{410}{mL/s} = \qty{25}{L/min}$.
\textbf{21.} $\num{24700}/180 = \qty{137}{mL}$.
\textbf{22.} $24.7\times(200 - 40) = \qty{3950}{mL/min}$, yaitu enam
belas kali \qty{250}{mL/min} ketika beristirahat.
\textbf{23.} Alirannya $\times 4.9$, penyaripatiannya $160/50 = \times 3.2$;
$4.9\times 3.2 = 16$.
\textbf{24.} Perintah dari korteks motoriknya, dan isyarat dari
reseptor ototnya, menimbang ulang masukan \omterm{def:b2:blood-pressure:baroreflex}{baroreseptornya} di dalam
medulanya sehingga refleksnya membela \qty{110}{mmHg} alih-alih 95 ---
titik acuannya dinaikkan, bukan gelungnya dilumpuhkan.
\textbf{25.} Penurunan sisa ketika berdiri \qty{5.7}{mmHg};
\qty{86}{mmHg} setelah pendarahannya; \qty{25}{L/min} dan \qty{4}{L}
oksigen tiap menit ketika berolahraga.
\end{solution}
